\documentclass[12pt,a4paper]{article}

\usepackage[hidelayout]{../../../Latex/Style/coursework-report-core}

\addbibresource{res-geo.bib}

\title{RES-GEO --- Barrier Geometry and Spatial Representation}
\author{}
\date{\today}

\begin{document}

\maketitle

\tableofcontents
\clearpage

\section{Domain Purpose and Scope}
\label{sec:res-geo-domain-purpose}

RES-GEO defines the geometry taxonomy, descriptors and software
representation needed for fixed-rigid barrier-class comparison in the
Kamaishi and Sendai local tsunami corridors. The active classes are
wall-type barriers and obstacle, dissipating or array-type barriers.

The current Research scope is not unrestricted geometry generation.
Existing or reconstructed defences may be included where source data
are sufficient, but the baseline comparison is between evidence-backed
barrier classes with explicit dimensions, placement rules and numerical
representation. ML-driven, surrogate-driven and generative geometry
work remains deferred until the hydrodynamic data pipeline and
validation hierarchy are accepted.

\section{Domain Deliverables}
\label{sec:res-geo-domain-deliverables}

% -------------------------------------------------------------------------
% WBS ID: RES-GEO-TAX
% Deliverable: Barrier Geometry Taxonomy and Candidate Families
% Status: In progress
% Evidence status: Barrier-class scope established; real-defence source gap open
% Review status: Not reviewed
% Last updated: 2026-07-19
% Dependencies: RES-DAT-SCN, RES-DAT-CAS, RES-PHY-BAR
% Downstream interfaces: RES-GEO-PAR, RES-GEO-MET, RES-GEO-REP, Software geometry manifest
% -------------------------------------------------------------------------

\section{RES-GEO-TAX --- Barrier Geometry Taxonomy and Candidate Families}
\label{sec:res-geo-tax}

\subsection{Purpose and Scope}
\label{sec:res-geo-tax-purpose}

This Deliverable classifies the active fixed-rigid barrier families
used for the Tohoku corridor comparison. It admits two baseline
families:

\begin{enumerate}
  \item wall-type barriers, including continuous or segmented seawalls,
    ridges, breakwaters and raised impermeable lines;
  \item obstacle, dissipating or array-type barriers, including finite
    blocks, columns, staggered rows, permeable-looking resolved
    openings and other structures whose effect is controlled by
    blockage, spacing and wake interaction.
\end{enumerate}

Selection is based on geometric recoverability, hydrodynamic relevance
and numerical representability. The deliverable excludes unrestricted
geometry generation and does not treat damage, scour or deformable
barrier response as active baseline geometry.

\subsection{Research Questions}
\label{sec:res-geo-tax-research-questions}

\begin{enumerate}
  \item Which barrier families are active in the Kamaishi and Sendai
    comparison cases?
  \item Which minimum descriptors distinguish wall-type and
    obstacle/array-type barriers for the local URANS--VOF model?
  \item Which real-defence geometries have enough provenance to be
    reconstructed rather than treated as TODO candidates?
\end{enumerate}

\subsection{Terminology and Definitions}
\label{sec:res-geo-tax-definitions}

\begin{description}
  \item[Wall-type barrier] A laterally extended defence whose principal
    hydrodynamic effect is reflection, run-up, overtopping and
    blocking along a near-continuous crest.
  \item[Obstacle or dissipating barrier] A finite-width body or row of
    bodies whose principal hydrodynamic effect is local blockage,
    separation, wake production, transmission and turbulent energy
    loss.
  \item[Array] A repeated set of resolved obstacle elements with
    spacing, staggering and row count recorded explicitly.
  \item[Fixed rigid barrier] A solid geometry that supplies no-slip and
    no-penetration wall patches to the local fluid solver and does not
    deform during the baseline hydrodynamic calculation.
\end{description}

\subsection{Literature Evidence}
\label{sec:res-geo-tax-literature-evidence}

\textcite{deljesus2012porous} supports the taxonomy split between
impermeable vertical barriers and barriers whose porosity or openings
require explicit modelling assumptions. Source role: supports the
wall-type versus dissipating/opening class distinction and limits any
claim about unresolved porous resistance. Project conclusion:
porosity-like behaviour is active only when openings are resolved by
geometry or a later porous closure is formally activated.

\textcite{watanabe2022validation} supports using controlled idealised
coastal structures for local hydrodynamic validation. Source role:
supports simple fixed obstacle and wall geometries as validation-facing
families rather than arbitrary generated forms. Limitation: the paper
does not define the final Kamaishi or Sendai defence geometry.

\textcite{mori2012nationwide} supports retaining Sanriku/Kamaishi and
Sendai coastal-plain settings as contrasting Tohoku impact regimes.
Source role: motivates site-facing geometry classes, not detailed
barrier dimensions. TODO/missing-source: identify a citable primary
source for the Kamaishi bay-mouth breakwater geometry, 2011 damage
state and any reconstruction dimensions before recreating it as a
project barrier.

\subsection{Governing Relations and Quantitative Definitions}
\label{sec:res-geo-tax-governing-relations}

For software-facing comparison, each geometry record shall include a
class label and a minimum descriptor vector:

\begin{align}
  \mathcal{G}_j
  &=
  \left\{
    c_j,\,
    \mathbf{x}_j,\,
    \psi_j,\,
    H_j,\,
    B_j,\,
    L_j,\,
    C_j,\,
    S_j,\,
    N_j,\,
    \chi_j
  \right\}
  \label{eq:res-geo-tax-minimum-geometry-record}
\end{align}

where $c_j$ is the barrier class, $\mathbf{x}_j$ is the placement
reference point in the local corridor coordinate system, $\psi_j$ is
the plan-view orientation, $H_j$ is height, $B_j$ is width or
thickness, $L_j$ is along-barrier length, $C_j$ is crest elevation,
$S_j$ is spacing for arrays, $N_j$ is element or row count and
$\chi_j$ records resolved openings, roughness and provenance flags.
Decision role: \cref{eq:res-geo-tax-minimum-geometry-record} is a data
contract, not a design-generation equation. Supporting evidence:
\textcite{deljesus2012porous} and \textcite{watanabe2022validation}.
Limitation: final admissible parameter bounds remain site- and
mesh-dependent.

\subsection{Geometry Family or Representation}
\label{sec:res-geo-tax-geometry-family-representation}

The active family set is:

\begin{itemize}
  \item \texttt{wall}: continuous or segmented impermeable line
    defence, with optional resolved openings;
  \item \texttt{obstacle}: isolated block, pier, mound or short
    barrier element;
  \item \texttt{array}: repeated obstacle elements whose spacing,
    staggering and row count are explicit;
  \item \texttt{dissipating}: resolved obstacle or array geometry
    intended to increase separation, wake loss, reflection or
    overtopping reduction.
\end{itemize}

Unresolved porous-media barriers, compliant structures, scour-adaptive
foundations and damage-evolving barriers are deferred extensions.

\subsection{Design Variables and Parameter Bounds}
\label{sec:res-geo-tax-design-variables-parameter-bounds}

% TODO: Complete this RES-GEO Domain-specific subsection.

\subsection{Geometric Descriptors}
\label{sec:res-geo-tax-geometric-descriptors}

% TODO: Complete this RES-GEO Domain-specific subsection.

\subsection{Placement and Arrangement Rules}
\label{sec:res-geo-tax-placement-arrangement-rules}

% TODO: Complete this RES-GEO Domain-specific subsection.

\subsection{Feasibility Constraints}
\label{sec:res-geo-tax-feasibility-constraints}

% TODO: Complete this RES-GEO Domain-specific subsection.

\subsection{Two- and Three-Dimensional Representation}
\label{sec:res-geo-tax-two-three-dimensional-representation}

% TODO: Complete this RES-GEO Domain-specific subsection.

\subsection{Candidate Approaches}
\label{sec:res-geo-tax-candidate-approaches}

% TODO: Define the credible candidate approaches considered.

\subsection{Critical Comparison}
\label{sec:res-geo-tax-critical-comparison}

% TODO: Compare candidates using physically and project-relevant criteria.
% Do not select an approach solely on popularity or implementation ease.

\subsection{Assumptions and Applicability}
\label{sec:res-geo-tax-assumptions}

% TODO: State assumptions and the conditions under which they are valid.

\subsection{Limitations and Failure Conditions}
\label{sec:res-geo-tax-limitations}

% TODO: State where the evidence, model or selected approach ceases to be
% reliable or applicable.

\subsection{Project-Specific Conclusions}
\label{sec:res-geo-tax-conclusions}

The project taxonomy is now class-comparison led. A geometry may enter
the active campaign when it has a traceable descriptor record, a fixed
rigid local representation and a clear hydrodynamic purpose. Real
defence recreation is allowed only after source geometry is identified;
otherwise the case remains a TODO rather than an invented citation.

\subsection{Selected Approach and Rationale}
\label{sec:res-geo-tax-selected-approach}

The selected approach is a small active taxonomy: wall-type barriers
and obstacle/dissipating/array barriers. The rejected active approach
is unrestricted generated geometry, because the present work must first
stabilise data handling, local hydrodynamic load extraction and the
V1--V2--V3 validation hierarchy.

\subsection{Unresolved Decisions}
\label{sec:res-geo-tax-unresolved-decisions}

\begin{itemize}
  \item TODO: obtain a primary or peer-reviewed source for Kamaishi
    bay-mouth breakwater dimensions and 2011 condition.
  \item Select final descriptor bounds after RES-DAT-GEO terrain
    extraction and RES-NUM-MSH mesh constraints are available.
  \item Decide whether any dissipating class requires a porous closure
    rather than resolved openings.
\end{itemize}

\subsection{Requirements and Handoffs}
\label{sec:res-geo-tax-requirements}

Software shall support a geometry manifest containing the descriptor
record in \cref{eq:res-geo-tax-minimum-geometry-record}, class labels,
local corridor coordinates, source/provenance fields and flags for
resolved openings versus deferred porous physics. `RES-MOD-IBC` and
`RES-NUM-MSH` shall consume this record as fixed terrain/barrier wall
patches; `RES-VER-MET` shall use the same identifiers for load and
energy metrics.

\subsection{Deliverable Completion Record}
\label{sec:res-geo-tax-completion-record}

% TODO: Record whether the Deliverable is:
%
% - Not started
% - In progress
% - Draft complete
% - Reviewed
% - Accepted
%
% Acceptance requires:
%
% - the research questions to be answered;
% - claims to be supported by appropriate evidence;
% - assumptions and limitations to be explicit;
% - the project-specific conclusion to be stated;
% - downstream requirements to be traceable;
% - unresolved decisions to be closed or formally deferred.

% -------------------------------------------------------------------------
% WBS ID: RES-GEO-PAR
% Deliverable: Parametric Geometry and Design-Variable Definition
% Status: In progress
% Evidence status: Active comparison descriptor set established
% Review status: Not reviewed
% Last updated: 2026-07-19
% Dependencies: RES-GEO-TAX, RES-DAT-SCN, RES-NUM-MSH
% Downstream interfaces: RES-GEO-MET, RES-GEO-REP, Software geometry manifest
% -------------------------------------------------------------------------

\section{RES-GEO-PAR --- Parametric Geometry and Design-Variable Definition}
\label{sec:res-geo-par}

\subsection{Purpose and Scope}
\label{sec:res-geo-par-purpose}

This Deliverable defines the bounded parameters required to compare
fixed-rigid barrier classes in the two local corridors. It does not
define an optimisation-ready free-form design space and does not
activate ML or generative geometry.

\subsection{Research Questions}
\label{sec:res-geo-par-research-questions}

\begin{enumerate}
  \item Which geometric parameters are required for wall-type and
    obstacle/array-type barriers?
  \item Which parameters are measured, inferred, scenario-controlled or
    unresolved?
  \item Which parameter bounds are controlled by terrain coverage,
    corridor width and meshability rather than design preference?
\end{enumerate}

\subsection{Terminology and Definitions}
\label{sec:res-geo-par-definitions}

% TODO: Define the technical terminology, symbols, quantities and
% classifications used in this Deliverable.

\subsection{Literature Evidence}
\label{sec:res-geo-par-literature-evidence}

\textcite{deljesus2012porous} supports separate parameter treatment
for impermeable walls and barriers whose openings or porosity change
hydraulic response. Source role: supports explicit opening/spacing
parameters and limits unresolved porous claims. \textcite{watanabe2022validation}
supports simple controlled geometry parameters for local impact and
loading validation. Project conclusion: the parameter set shall remain
small, inspectable and tied to hydrodynamic metrics.

\subsection{Governing Relations and Quantitative Definitions}
\label{sec:res-geo-par-governing-relations}

For each case $j$, the comparison parameter vector is:

\begin{align}
  \boldsymbol{\theta}^{\mathrm{geo}}_j
  &=
  \left[
    H_j,\,
    C_j,\,
    B_j,\,
    L_j,\,
    \psi_j,\,
    x'_j,\,
    y'_j,\,
    S_j,\,
    N_j,\,
    \phi^{\mathrm{open}}_j
  \right]^{\mathsf{T}}
  \label{eq:res-geo-par-parameter-vector}
\end{align}

where $\phi^{\mathrm{open}}_j$ is the resolved open-area or gap
fraction for array or segmented cases. Source role:
\textcite{deljesus2012porous} supports separating impermeable and
opening-controlled barriers; \textcite{watanabe2022validation}
supports maintaining geometry parameters that map directly to local
loading observables. Limitation: \cref{eq:res-geo-par-parameter-vector}
does not set final bounds.

\subsection{Geometry Family or Representation}
\label{sec:res-geo-par-geometry-family-representation}

% TODO: Complete this RES-GEO Domain-specific subsection.

\subsection{Design Variables and Parameter Bounds}
\label{sec:res-geo-par-design-variables-parameter-bounds}

Parameter sets shall distinguish measured geometry, inferred geometry,
scenario-controlled geometry and deferred structural assumptions.
Candidate parameters may include crest height, barrier height, wall or
element thickness, footprint length, plan orientation, corridor
placement, gap spacing, row count, resolved opening fraction and
surface roughness flag. Foundation strengthening, material changes,
damage evolution and deformable response are stretch extensions owned
by RES-STR after hydrodynamic load extraction is stable.

\subsection{Geometric Descriptors}
\label{sec:res-geo-par-geometric-descriptors}

% TODO: Complete this RES-GEO Domain-specific subsection.

\subsection{Placement and Arrangement Rules}
\label{sec:res-geo-par-placement-arrangement-rules}

% TODO: Complete this RES-GEO Domain-specific subsection.

\subsection{Feasibility Constraints}
\label{sec:res-geo-par-feasibility-constraints}

% TODO: Complete this RES-GEO Domain-specific subsection.

\subsection{Two- and Three-Dimensional Representation}
\label{sec:res-geo-par-two-three-dimensional-representation}

% TODO: Complete this RES-GEO Domain-specific subsection.

\subsection{Candidate Approaches}
\label{sec:res-geo-par-candidate-approaches}

% TODO: Define the credible candidate approaches considered.

\subsection{Critical Comparison}
\label{sec:res-geo-par-critical-comparison}

% TODO: Compare candidates using physically and project-relevant criteria.
% Do not select an approach solely on popularity or implementation ease.

\subsection{Assumptions and Applicability}
\label{sec:res-geo-par-assumptions}

% TODO: State assumptions and the conditions under which they are valid.

\subsection{Limitations and Failure Conditions}
\label{sec:res-geo-par-limitations}

% TODO: State where the evidence, model or selected approach ceases to be
% reliable or applicable.

\subsection{Project-Specific Conclusions}
\label{sec:res-geo-par-conclusions}

The active Studentship comparison requires a compact parameter vector,
not an open-ended design grammar. Geometry choices must be reproducible
from the manifest, meshable in the local domain and traceable to
barrier-performance metrics.

\subsection{Selected Approach and Rationale}
\label{sec:res-geo-par-selected-approach}

The selected approach is a descriptor-first parameter set for fixed
wall, obstacle, dissipating and array barriers. Broad geometric
optimisation and generative parameter discovery are formally deferred.

\subsection{Unresolved Decisions}
\label{sec:res-geo-par-unresolved-decisions}

\begin{itemize}
  \item Final numeric bounds for $H_j$, $B_j$, $L_j$, $S_j$ and
    $N_j$ remain open until corridor terrain and mesh constraints are
    accepted.
  \item Real-defence reconstruction parameters for Kamaishi remain
    TODO pending a citable geometry source.
  \item Roughness and opening/porosity flags remain provisional until
    RES-MOD-SRC decides whether any porous closure is active.
\end{itemize}

\subsection{Requirements and Handoffs}
\label{sec:res-geo-par-requirements}

Software shall ingest \cref{eq:res-geo-par-parameter-vector} from the
geometry manifest, generate fixed barrier patches, and preserve stable
case identifiers for metric extraction. RES-NUM-MSH shall report
whether each parameter set is meshable; RES-VER-MET shall record which
geometry parameters produced each run-up, overtopping, force, moment
and energy metric.

\subsection{Deliverable Completion Record}
\label{sec:res-geo-par-completion-record}

% TODO: Record whether the Deliverable is:
%
% - Not started
% - In progress
% - Draft complete
% - Reviewed
% - Accepted
%
% Acceptance requires:
%
% - the research questions to be answered;
% - claims to be supported by appropriate evidence;
% - assumptions and limitations to be explicit;
% - the project-specific conclusion to be stated;
% - downstream requirements to be traceable;
% - unresolved decisions to be closed or formally deferred.

% -------------------------------------------------------------------------
% WBS ID: RES-GEO-MET
% Deliverable: Geometric and Hydraulic Descriptor Framework
% Status: In progress
% Evidence status: Minimum descriptors tied to hydrodynamic outputs
% Review status: Not reviewed
% Last updated: 2026-07-19
% Dependencies: RES-GEO-PAR, RES-DAT-SCN
% Downstream interfaces: RES-PHY-BAR, RES-PHY-LOD, RES-VER-MET, Software output schema
% -------------------------------------------------------------------------

\section{RES-GEO-MET --- Geometric and Hydraulic Descriptor Framework}
\label{sec:res-geo-met}

\subsection{Purpose and Scope}
\label{sec:res-geo-met-purpose}

This Deliverable defines the minimum geometric and hydraulic
descriptors needed to compare wall-type and obstacle/array-type
barriers in fixed-rigid local tsunami simulations. Descriptor records
shall preserve provenance and uncertainty. Damage state, material
deformation and scour descriptors are recorded only as deferred
structural extensions.

\subsection{Research Questions}
\label{sec:res-geo-met-research-questions}

\begin{enumerate}
  \item Which descriptors are required before local hydrodynamic
    comparison can begin?
  \item Which descriptors drive software geometry, and which are
    post-processing metrics emitted by the solver?
  \item Which descriptors must carry uncertainty/provenance flags?
\end{enumerate}

\subsection{Terminology and Definitions}
\label{sec:res-geo-met-definitions}

% TODO: Define the technical terminology, symbols, quantities and
% classifications used in this Deliverable.

\subsection{Literature Evidence}
\label{sec:res-geo-met-literature-evidence}

\textcite{deljesus2012porous} links vertical barrier form, porosity
and transmission to hydraulic response. Source role: supports
including blockage, opening and transmission descriptors.
\textcite{watanabe2022validation} links local geometry to separate
force, pressure, velocity and free-surface validation quantities.
Project conclusion: descriptors must identify both what was simulated
and which hydrodynamic outputs are meaningful for comparison.

\subsection{Governing Relations and Quantitative Definitions}
\label{sec:res-geo-met-governing-relations}

For array and obstacle cases, the plan blockage ratio shall be
reported as:

\begin{align}
  \beta_{\mathrm{blk}}
  &=
  \frac{
    A_{\perp,\mathrm{solid}}
  }{
    A_{\perp,\mathrm{corridor}}
    +
    \epsilon_A
  }
  \label{eq:res-geo-met-blockage-ratio}
\end{align}

where $A_{\perp,\mathrm{solid}}$ is the projected solid area normal to
the corridor centreline and $A_{\perp,\mathrm{corridor}}$ is the
available projected corridor area at the comparison section. Decision
role: \cref{eq:res-geo-met-blockage-ratio} is a comparison descriptor
for wall/obstacle/array cases. Supporting evidence:
\textcite{deljesus2012porous}. Limitation: it does not replace full
3D flow solution because wake, overtopping and pressure depend on
local geometry and bathymetry.

\subsection{Geometry Family or Representation}
\label{sec:res-geo-met-geometry-family-representation}

% TODO: Complete this RES-GEO Domain-specific subsection.

\subsection{Design Variables and Parameter Bounds}
\label{sec:res-geo-met-design-variables-parameter-bounds}

% TODO: Complete this RES-GEO Domain-specific subsection.

\subsection{Geometric Descriptors}
\label{sec:res-geo-met-geometric-descriptors}

Required pre-simulation descriptors are class label, local corridor
coordinates, orientation, crest elevation, height, thickness or
element width, footprint length, spacing, row count, opening fraction,
roughness flag, blockage ratio and provenance quality.

Required post-simulation hydraulic descriptors are maximum run-up,
overtopping volume or discharge, transmitted and reflected energy
measures, downstream maximum depth and velocity, peak pressure,
resultant force, impulse and moment. Damage, survival state and
structural-capacity proxies are retained only for later RES-STR
screening and shall not be used to claim baseline hydrodynamic
performance.

\subsection{Placement and Arrangement Rules}
\label{sec:res-geo-met-placement-arrangement-rules}

% TODO: Complete this RES-GEO Domain-specific subsection.

\subsection{Feasibility Constraints}
\label{sec:res-geo-met-feasibility-constraints}

% TODO: Complete this RES-GEO Domain-specific subsection.

\subsection{Two- and Three-Dimensional Representation}
\label{sec:res-geo-met-two-three-dimensional-representation}

% TODO: Complete this RES-GEO Domain-specific subsection.

\subsection{Candidate Approaches}
\label{sec:res-geo-met-candidate-approaches}

% TODO: Define the credible candidate approaches considered.

\subsection{Critical Comparison}
\label{sec:res-geo-met-critical-comparison}

% TODO: Compare candidates using physically and project-relevant criteria.
% Do not select an approach solely on popularity or implementation ease.

\subsection{Assumptions and Applicability}
\label{sec:res-geo-met-assumptions}

% TODO: State assumptions and the conditions under which they are valid.

\subsection{Limitations and Failure Conditions}
\label{sec:res-geo-met-limitations}

% TODO: State where the evidence, model or selected approach ceases to be
% reliable or applicable.

\subsection{Project-Specific Conclusions}
\label{sec:res-geo-met-conclusions}

The geometry descriptor framework is sufficient when each barrier case
can be reconstructed, meshed, run and post-processed using stable
identifiers. It is not sufficient for structural certification or
damage prediction.

\subsection{Selected Approach and Rationale}
\label{sec:res-geo-met-selected-approach}

The selected descriptor approach separates pre-simulation geometry
fields from solver-emitted hydrodynamic metrics. This keeps software
data handling parallel with Research without moving solver
implementation into RES-GEO.

\subsection{Unresolved Decisions}
\label{sec:res-geo-met-unresolved-decisions}

Open decisions are the final blockage-normalisation section, roughness
descriptor values, gap/opening representation and accepted uncertainty
classes for inferred geometry.

\subsection{Requirements and Handoffs}
\label{sec:res-geo-met-requirements}

Software shall emit descriptors and metrics using the same
\texttt{corridor\_id}, \texttt{geometry\_id} and \texttt{case\_id}.
RES-PHY-BAR and RES-PHY-LOD shall define the hydrodynamic meaning of
the emitted metrics; RES-VER-MET shall define acceptance tolerances.

\subsection{Deliverable Completion Record}
\label{sec:res-geo-met-completion-record}

% TODO: Record whether the Deliverable is:
%
% - Not started
% - In progress
% - Draft complete
% - Reviewed
% - Accepted
%
% Acceptance requires:
%
% - the research questions to be answered;
% - claims to be supported by appropriate evidence;
% - assumptions and limitations to be explicit;
% - the project-specific conclusion to be stated;
% - downstream requirements to be traceable;
% - unresolved decisions to be closed or formally deferred.

% -------------------------------------------------------------------------
% WBS ID: RES-GEO-ARR
% Deliverable: Barrier Arrangement, Placement and Spatial Configuration
% Status: In progress
% Evidence status: Corridor-linked placement rules established
% Review status: Not reviewed
% Last updated: 2026-07-19
% Dependencies: RES-DAT-SCN, RES-GEO-TAX, RES-GEO-PAR
% Downstream interfaces: RES-GEO-REP, RES-MOD-IBC, RES-NUM-MSH, Software geometry manifest
% -------------------------------------------------------------------------

\section{RES-GEO-ARR --- Barrier Arrangement, Placement and Spatial Configuration}
\label{sec:res-geo-arr}

\subsection{Purpose and Scope}
\label{sec:res-geo-arr-purpose}

This Deliverable defines how active barrier classes are placed inside
the local corridor coordinate system. Arrangement is constrained by
the evidence-based Kamaishi and Sendai corridor polygons, not by an
arbitrary drawing canvas. The deliverable excludes global coastline
redesign and generative placement optimisation.

\subsection{Research Questions}
\label{sec:res-geo-arr-research-questions}

\begin{enumerate}
  \item How is a barrier positioned relative to the corridor
    centreline $x'$ and width direction $y'$?
  \item How are continuous walls distinguished from finite obstacles
    and arrays?
  \item Which placement choices are physical scenario decisions, and
    which are numerical buffer or mesh-quality decisions?
\end{enumerate}

\subsection{Terminology and Definitions}
\label{sec:res-geo-arr-definitions}

$x'$ is the local propagation-direction coordinate supplied by
RES-DAT-SCN, $y'$ is the transverse corridor coordinate, and $z'$ is
vertical elevation in the solver datum. The angle $\psi_j$ is the
plan-view orientation of barrier $j$ measured relative to the
$x'$ axis; $\psi_j=\pi/2$ denotes a barrier aligned across the
corridor.

\subsection{Literature Evidence}
\label{sec:res-geo-arr-literature-evidence}

\textcite{mori2012nationwide} supports retaining different placement
contexts for Sanriku/Kamaishi and Sendai because the same event
produced contrasting run-up and inundation behaviour. Source role:
supports corridor-specific placement rather than one generic test
site. \textcite{watanabe2022validation} supports preserving exact
probe and structure locations for validation-facing local simulations.
Limitation: neither source supplies final barrier coordinates; those
remain RES-DAT-SCN/RES-DAT-CAS data products.

\subsection{Governing Relations and Quantitative Definitions}
\label{sec:res-geo-arr-governing-relations}

The plan-view local placement vector is:

\begin{align}
  \mathbf{r}'_j
  &=
  \left[
    x'_j,\,
    y'_j
  \right]^{\mathsf{T}}
  \label{eq:res-geo-arr-placement-vector}
\end{align}

The local basis for each barrier is:

\begin{align}
  \mathbf{e}_{n,j}
  &=
  \left[
    \cos\psi_j,\,
    \sin\psi_j
  \right]^{\mathsf{T}}
  \label{eq:res-geo-arr-barrier-normal}
  \\
  \mathbf{e}_{t,j}
  &=
  \left[
    -\sin\psi_j,\,
    \cos\psi_j
  \right]^{\mathsf{T}}
  \label{eq:res-geo-arr-barrier-tangent}
\end{align}

Decision role: these definitions make barrier placement reproducible
from the corridor transform. Supporting evidence: RES-DAT-SCN owns
the corridor construction; \textcite{watanabe2022validation} supports
validation-facing control of structure and probe locations.

\subsection{Geometry Family or Representation}
\label{sec:res-geo-arr-geometry-family-representation}

Continuous walls are represented as elongated patches whose tangent
direction is $\mathbf{e}_{t,j}$. Obstacle and array cases are
represented as one or more finite elements with recorded centre
points, spacing and row index.

\subsection{Design Variables and Parameter Bounds}
\label{sec:res-geo-arr-design-variables-parameter-bounds}

Arrangement variables include $x'_j$, $y'_j$, $\psi_j$, number of rows,
along-row spacing, cross-row spacing, stagger offset and gap/opening
fraction. These variables are comparison controls, not yet an
optimisation search space.

\subsection{Geometric Descriptors}
\label{sec:res-geo-arr-geometric-descriptors}

Placement descriptors shall include clearance to artificial inlet,
outlet and side boundaries; clearance to steep terrain changes; and
distance from run-up, overtopping and wake measurement sections.

\subsection{Placement and Arrangement Rules}
\label{sec:res-geo-arr-placement-arrangement-rules}

Baseline placement rules are:

\begin{enumerate}
  \item place wall-type cases approximately transverse to $x'$ unless
    a real coastline, harbour or defence alignment provides a citable
    alternative;
  \item place array-type cases with all elements inside the accepted
    terrain/corridor polygon and away from sponge zones;
  \item keep comparison sections fixed when comparing barrier classes;
  \item do not narrow the corridor to fit a barrier unless the
    narrowing represents a physical bay or coastline constraint.
\end{enumerate}

\subsection{Feasibility Constraints}
\label{sec:res-geo-arr-feasibility-constraints}

A placement is infeasible if it intersects missing terrain data,
requires a nonphysical wall at the side boundary, violates local mesh
quality constraints or places the barrier inside a numerical
relaxation region.

\subsection{Two- and Three-Dimensional Representation}
\label{sec:res-geo-arr-two-three-dimensional-representation}

% TODO: Complete this RES-GEO Domain-specific subsection.

\subsection{Candidate Approaches}
\label{sec:res-geo-arr-candidate-approaches}

% TODO: Define the credible candidate approaches considered.

\subsection{Critical Comparison}
\label{sec:res-geo-arr-critical-comparison}

% TODO: Compare candidates using physically and project-relevant criteria.
% Do not select an approach solely on popularity or implementation ease.

\subsection{Assumptions and Applicability}
\label{sec:res-geo-arr-assumptions}

% TODO: State assumptions and the conditions under which they are valid.

\subsection{Limitations and Failure Conditions}
\label{sec:res-geo-arr-limitations}

% TODO: State where the evidence, model or selected approach ceases to be
% reliable or applicable.

\subsection{Project-Specific Conclusions}
\label{sec:res-geo-arr-conclusions}

For the Studentship campaign, arrangement is part of the data/model
interface. It links source-to-coast trajectory, terrain extraction,
barrier geometry and validation probes; it is not a separate design
art exercise.

\subsection{Selected Approach and Rationale}
\label{sec:res-geo-arr-selected-approach}

The selected approach is corridor-relative placement using
$(x',y',z')$ coordinates and fixed comparison sections. Free placement
optimisation is deferred.

\subsection{Unresolved Decisions}
\label{sec:res-geo-arr-unresolved-decisions}

Open decisions are final Kamaishi and Sendai target coordinates, final
comparison sections, sponge widths, and any real-defence alignment
source for the Kamaishi breakwater.

\subsection{Requirements and Handoffs}
\label{sec:res-geo-arr-requirements}

Software shall generate barrier placement from corridor coordinates,
clip it against accepted terrain/corridor polygons, and emit placement
diagnostics for mesh quality and validation probe locations.

\subsection{Deliverable Completion Record}
\label{sec:res-geo-arr-completion-record}

% TODO: Record whether the Deliverable is:
%
% - Not started
% - In progress
% - Draft complete
% - Reviewed
% - Accepted
%
% Acceptance requires:
%
% - the research questions to be answered;
% - claims to be supported by appropriate evidence;
% - assumptions and limitations to be explicit;
% - the project-specific conclusion to be stated;
% - downstream requirements to be traceable;
% - unresolved decisions to be closed or formally deferred.

% -------------------------------------------------------------------------
% WBS ID: RES-GEO-REP
% Deliverable: Two- and Three-Dimensional Model Representation
% Status: In progress
% Evidence status: Fixed-rigid local representation established
% Review status: Not reviewed
% Last updated: 2026-07-19
% Dependencies: RES-GEO-PAR, RES-GEO-MET, RES-NUM-SPEC
% Downstream interfaces: RES-MOD-IBC, RES-NUM-MSH, RES-STR-LOD, RES-VER-SPEC, Software mesh/field containers
% -------------------------------------------------------------------------

\section{RES-GEO-REP --- Two- and Three-Dimensional Model Representation}
\label{sec:res-geo-rep}

\subsection{Purpose and Scope}
\label{sec:res-geo-rep-purpose}

This Deliverable defines how each active fixed-rigid barrier geometry
is represented for the regional model, local 3D fluid model,
hydrodynamic-load handoff and later deferred structural analysis. ML
feature generation is not active; any feature representation is a
future consumer of the same manifest.

\subsection{Research Questions}
\label{sec:res-geo-rep-research-questions}

\begin{enumerate}
  \item What is the minimum regional footprint required for regional
    propagation and replay extraction?
  \item What local 3D surface representation is required for
    no-slip/no-penetration wall patches and load extraction?
  \item Which representation fields must be preserved for later
    structural or ML use without activating those workstreams now?
\end{enumerate}

\subsection{Terminology and Definitions}
\label{sec:res-geo-rep-definitions}

% TODO: Define the technical terminology, symbols, quantities and
% classifications used in this Deliverable.

\subsection{Literature Evidence}
\label{sec:res-geo-rep-literature-evidence}

\textcite{watanabe2022validation} supports explicit local geometry,
probe and loading-output control in tsunami-impact validation.
\textcite{deljesus2012porous} supports distinguishing resolved solid
geometry from unresolved porous closure. Project conclusion: the
local representation shall be a resolved fixed solid surface unless a
separate porous model is later accepted.

\subsection{Governing Relations and Quantitative Definitions}
\label{sec:res-geo-rep-governing-relations}

The fixed barrier patch in the local fluid domain is:

\begin{align}
  \Gamma_{\mathrm{B}}
  &=
  \partial\Omega_{\mathrm{f}}
  \cap
  \partial\Omega_{\mathrm{B}}
  \label{eq:res-geo-rep-fixed-barrier-patch}
\end{align}

where $\Omega_{\mathrm{f}}$ is the local fluid-control domain and
$\Omega_{\mathrm{B}}$ is the solid barrier volume. Decision role:
\cref{eq:res-geo-rep-fixed-barrier-patch} defines the geometric
surface on which no-slip/no-penetration conditions and hydrodynamic
load integration are applied. Supporting evidence:
\textcite{watanabe2022validation}. Limitation: it does not describe
deformable FSI interfaces.

\subsection{Geometry Family or Representation}
\label{sec:res-geo-rep-geometry-family-representation}

% TODO: Complete this RES-GEO Domain-specific subsection.

\subsection{Design Variables and Parameter Bounds}
\label{sec:res-geo-rep-design-variables-parameter-bounds}

% TODO: Complete this RES-GEO Domain-specific subsection.

\subsection{Geometric Descriptors}
\label{sec:res-geo-rep-geometric-descriptors}

% TODO: Complete this RES-GEO Domain-specific subsection.

\subsection{Placement and Arrangement Rules}
\label{sec:res-geo-rep-placement-arrangement-rules}

% TODO: Complete this RES-GEO Domain-specific subsection.

\subsection{Feasibility Constraints}
\label{sec:res-geo-rep-feasibility-constraints}

% TODO: Complete this RES-GEO Domain-specific subsection.

\subsection{Two- and Three-Dimensional Representation}
\label{sec:res-geo-rep-two-three-dimensional-representation}

Representation levels shall include:

\begin{itemize}
  \item a regional-scale footprint or roughness/obstruction marker,
    used only if the regional model needs a defence representation;
  \item a local three-dimensional fixed solid surface for the
    URANS--VOF barrier interaction model;
  \item a patch and face identifier set for pressure, shear, force,
    impulse and moment extraction;
  \item a deferred structural surface mapping for one-way CFD-to-FEA
    load transfer after hydrodynamic extraction is stable;
  \item optional future ML descriptors generated from the same
    manifest, not an independent geometry source.
\end{itemize}

\subsection{Candidate Approaches}
\label{sec:res-geo-rep-candidate-approaches}

% TODO: Define the credible candidate approaches considered.

\subsection{Critical Comparison}
\label{sec:res-geo-rep-critical-comparison}

% TODO: Compare candidates using physically and project-relevant criteria.
% Do not select an approach solely on popularity or implementation ease.

\subsection{Assumptions and Applicability}
\label{sec:res-geo-rep-assumptions}

% TODO: State assumptions and the conditions under which they are valid.

\subsection{Limitations and Failure Conditions}
\label{sec:res-geo-rep-limitations}

% TODO: State where the evidence, model or selected approach ceases to be
% reliable or applicable.

\subsection{Project-Specific Conclusions}
\label{sec:res-geo-rep-conclusions}

The local solver needs an explicit fixed surface representation before
any barrier comparison is meaningful. The same surface identifiers
must survive into post-processing so forces and moments can be traced
back to geometry cases.

\subsection{Selected Approach and Rationale}
\label{sec:res-geo-rep-selected-approach}

The selected local representation is a fixed rigid solid volume
converted to terrain/barrier wall patches and stable patch identifiers.
Moving boundaries, topology change and deformable surfaces are
deferred to RES-STR-FSI.

\subsection{Unresolved Decisions}
\label{sec:res-geo-rep-unresolved-decisions}

Open decisions are the exact geometry file format, wall roughness
metadata, CAD/CSG versus raster-derived surface generation, and
Kamaishi real-defence source geometry.

\subsection{Requirements and Handoffs}
\label{sec:res-geo-rep-requirements}

Software shall support fixed geometry ingestion, patch labelling,
face/patch metric extraction and structured export to HDF5 or an
equivalent schema. RES-STR-LOD shall consume the same patch identifiers
for later hydrodynamic load transfer.

\subsection{Deliverable Completion Record}
\label{sec:res-geo-rep-completion-record}

% TODO: Record whether the Deliverable is:
%
% - Not started
% - In progress
% - Draft complete
% - Reviewed
% - Accepted
%
% Acceptance requires:
%
% - the research questions to be answered;
% - claims to be supported by appropriate evidence;
% - assumptions and limitations to be explicit;
% - the project-specific conclusion to be stated;
% - downstream requirements to be traceable;
% - unresolved decisions to be closed or formally deferred.

% -------------------------------------------------------------------------
% WBS ID: RES-GEO-CST
% Deliverable: Geometric Feasibility and Design Constraints
% Status: In progress
% Evidence status: Fixed-rigid geometry feasibility gates established
% Review status: Not reviewed
% Last updated: 2026-07-19
% Dependencies: RES-GEO-PAR, RES-DAT-GEO, RES-NUM-MSH
% Downstream interfaces: RES-GEO-SPEC, RES-NUM-MSH, Software geometry manifest
% -------------------------------------------------------------------------

\section{RES-GEO-CST --- Geometric Feasibility and Design Constraints}
\label{sec:res-geo-cst}

\subsection{Purpose and Scope}
\label{sec:res-geo-cst-purpose}

This Deliverable narrows admissible geometry cases to fixed-rigid
barrier-class comparisons that can be represented on the accepted
terrain, meshed robustly and interpreted through hydrodynamic outputs.
Material limits, foundation constraints, deformability and damage are
recorded as stretch structural concerns rather than active geometric
feasibility criteria.

\subsection{Research Questions}
\label{sec:res-geo-cst-research-questions}

% TODO: State the questions that must be answered before the Deliverable
% can be accepted.

\subsection{Terminology and Definitions}
\label{sec:res-geo-cst-definitions}

% TODO: Define the technical terminology, symbols, quantities and
% classifications used in this Deliverable.

\subsection{Literature Evidence}
\label{sec:res-geo-cst-literature-evidence}

% TODO: Record evidence from peer-reviewed papers, authoritative datasets,
% standards, technical reports and primary documentation.
%
% Do not create a detached literature-summary list. Explain how each source
% supports, contradicts or limits a project decision.

\subsection{Governing Relations and Quantitative Definitions}
\label{sec:res-geo-cst-governing-relations}

% TODO: Record relevant equations, dimensionless groups, empirical models,
% extraction rules or quantitative definitions.
%
% Use align environments for displayed equations.
% Every displayed equation must have a unique \label{eq:...}.
% Do not place punctuation after displayed equations.

\subsection{Geometry Family or Representation}
\label{sec:res-geo-cst-geometry-family-representation}

% TODO: Complete this RES-GEO Domain-specific subsection.

\subsection{Design Variables and Parameter Bounds}
\label{sec:res-geo-cst-design-variables-parameter-bounds}

% TODO: Complete this RES-GEO Domain-specific subsection.

\subsection{Geometric Descriptors}
\label{sec:res-geo-cst-geometric-descriptors}

% TODO: Complete this RES-GEO Domain-specific subsection.

\subsection{Placement and Arrangement Rules}
\label{sec:res-geo-cst-placement-arrangement-rules}

% TODO: Complete this RES-GEO Domain-specific subsection.

\subsection{Feasibility Constraints}
\label{sec:res-geo-cst-feasibility-constraints}

Candidate interventions shall be rejected or marked unresolved if they cannot
be constructed plausibly, cannot be represented in the selected numerical
workflow, violate known material or foundation constraints, or lack enough
provenance to support simulation.

For the active hydrodynamic baseline, a geometry shall also be rejected
or deferred if it requires unresolved porous resistance, moving
structural boundaries, scour-induced terrain change or material
damage to explain its intended performance.

\textbf{Selected approach.} Feasibility is first a data and numerical
representation gate: the case must fit within the accepted corridor
polygon, avoid sponge zones, preserve mesh quality near walls and have
a traceable descriptor record. Supporting evidence:
\textcite{watanabe2022validation} for validation-facing geometry
control and \textcite{deljesus2012porous} for separating impermeable
and porous/dissipating assumptions.

\subsection{Two- and Three-Dimensional Representation}
\label{sec:res-geo-cst-two-three-dimensional-representation}

% TODO: Complete this RES-GEO Domain-specific subsection.

\subsection{Candidate Approaches}
\label{sec:res-geo-cst-candidate-approaches}

% TODO: Define the credible candidate approaches considered.

\subsection{Critical Comparison}
\label{sec:res-geo-cst-critical-comparison}

% TODO: Compare candidates using physically and project-relevant criteria.
% Do not select an approach solely on popularity or implementation ease.

\subsection{Assumptions and Applicability}
\label{sec:res-geo-cst-assumptions}

% TODO: State assumptions and the conditions under which they are valid.

\subsection{Limitations and Failure Conditions}
\label{sec:res-geo-cst-limitations}

% TODO: State where the evidence, model or selected approach ceases to be
% reliable or applicable.

\subsection{Project-Specific Conclusions}
\label{sec:res-geo-cst-conclusions}

The Studentship shall not spend simulation effort on barrier shapes
whose geometry cannot be reproduced, clipped to the corridor and
post-processed for hydrodynamic metrics.

\subsection{Selected Approach and Rationale}
\label{sec:res-geo-cst-selected-approach}

% TODO: Record the selected approach only after sufficient evidence exists.
% State the alternatives rejected and the reasons for rejection.

\subsection{Unresolved Decisions}
\label{sec:res-geo-cst-unresolved-decisions}

Open decisions are final mesh-quality limits, minimum wall clearance
from artificial boundaries, and whether any dissipating concept is
admissible without a porous closure.

\subsection{Requirements and Handoffs}
\label{sec:res-geo-cst-requirements}

Software shall report a geometry-feasibility status before launching
local runs, including corridor intersection checks, minimum cell-size
estimates, wall clearance and unresolved-physics flags.

\subsection{Deliverable Completion Record}
\label{sec:res-geo-cst-completion-record}

% TODO: Record whether the Deliverable is:
%
% - Not started
% - In progress
% - Draft complete
% - Reviewed
% - Accepted
%
% Acceptance requires:
%
% - the research questions to be answered;
% - claims to be supported by appropriate evidence;
% - assumptions and limitations to be explicit;
% - the project-specific conclusion to be stated;
% - downstream requirements to be traceable;
% - unresolved decisions to be closed or formally deferred.

% -------------------------------------------------------------------------
% WBS ID: RES-GEO-SPEC
% Deliverable: Integrated Barrier-Class Geometry Specification
% Status: In progress
% Evidence status: Active barrier-class comparison scope established
% Review status: Not reviewed
% Last updated: 2026-07-19
% Dependencies: RES-GEO-TAX through RES-GEO-CST
% Downstream interfaces: RES-MOD-IBC, RES-NUM-MSH, RES-PHY-BAR, RES-PHY-LOD, RES-VER-MET, Software geometry manifest
% -------------------------------------------------------------------------

\section{RES-GEO-SPEC --- Integrated Barrier-Class Geometry Specification}
\label{sec:res-geo-spec}

\subsection{Purpose and Scope}
\label{sec:res-geo-spec-purpose}

This Deliverable consolidates the active fixed-rigid barrier geometry
specification. It preserves the WBS ID while changing the active output
away from geometry generation and toward reproducible wall-type and
obstacle/dissipating/array comparisons in the two Tohoku corridors.

\subsection{Research Questions}
\label{sec:res-geo-spec-research-questions}

\begin{enumerate}
  \item Are the active geometry classes, descriptors, placement rules
    and representations sufficient for local 3D hydrodynamic
    comparison?
  \item Which fields must the Software workstream implement in the
    geometry manifest before data handling can proceed?
  \item Which geometry ideas are explicitly deferred?
\end{enumerate}

\subsection{Terminology and Definitions}
\label{sec:res-geo-spec-definitions}

% TODO: Define the technical terminology, symbols, quantities and
% classifications used in this Deliverable.

\subsection{Literature Evidence}
\label{sec:res-geo-spec-literature-evidence}

\textcite{deljesus2012porous} supports the impermeable/opening-aware
barrier distinction and limits unresolved porous claims.
\textcite{watanabe2022validation} supports preserving geometry and
output identifiers for local impact validation. \textcite{mori2012nationwide}
supports retaining the two Tohoku corridor settings as contrasting
hydrodynamic contexts. Project conclusion: geometry specification is a
data and modelling contract for comparison, not a literature survey or
automatic design generator.

\subsection{Governing Relations and Quantitative Definitions}
\label{sec:res-geo-spec-governing-relations}

The minimum geometry-case contract is:

\begin{align}
  \mathcal{C}^{\mathrm{geo}}_j
  &=
  \left\{
    \mathtt{corridor\_id},\,
    \mathtt{geometry\_id},\,
    \mathcal{G}_j,\,
    \boldsymbol{\theta}^{\mathrm{geo}}_j,\,
    \Gamma_{\mathrm{B},j},\,
    \mathcal{P}_j
  \right\}
  \label{eq:res-geo-spec-case-contract}
\end{align}

where $\mathcal{G}_j$ is the descriptor record,
$\boldsymbol{\theta}^{\mathrm{geo}}_j$ is the parameter vector,
$\Gamma_{\mathrm{B},j}$ is the fixed barrier patch and
$\mathcal{P}_j$ records provenance, CRS/datum linkage and unresolved
physics flags. Source role: this equation integrates the decisions in
RES-GEO-TAX, RES-GEO-PAR and RES-GEO-REP. Limitation: it does not
define solver implementation or validation results.

\subsection{Geometry Family or Representation}
\label{sec:res-geo-spec-geometry-family-representation}

The active representation is fixed terrain plus fixed rigid barrier
patches. Wall-type and obstacle/array-type cases share the same
manifest structure so they can be compared through common
hydrodynamic metrics.

\subsection{Design Variables and Parameter Bounds}
\label{sec:res-geo-spec-design-variables-parameter-bounds}

Only descriptor-level comparison variables are active. Generative
control points, topology optimisation variables and learned geometry
latents are deferred.

\subsection{Geometric Descriptors}
\label{sec:res-geo-spec-geometric-descriptors}

Minimum descriptors are class, position, orientation, crest elevation,
height, width/thickness, footprint length, spacing, row count, opening
fraction, roughness flag, blockage ratio and provenance quality.

\subsection{Placement and Arrangement Rules}
\label{sec:res-geo-spec-placement-arrangement-rules}

Placement shall be corridor-relative and shall preserve fixed
measurement sections across class comparisons.

\subsection{Feasibility Constraints}
\label{sec:res-geo-spec-feasibility-constraints}

Feasibility requires terrain coverage, meshability, stable wall patch
labels and no unresolved physics dependency for the active baseline.

\subsection{Two- and Three-Dimensional Representation}
\label{sec:res-geo-spec-two-three-dimensional-representation}

The local three-dimensional representation is a fixed barrier surface
used for no-slip/no-penetration boundary conditions and for pressure,
shear, force, impulse and moment extraction.

\subsection{Candidate Approaches}
\label{sec:res-geo-spec-candidate-approaches}

% TODO: Define the credible candidate approaches considered.

\subsection{Critical Comparison}
\label{sec:res-geo-spec-critical-comparison}

% TODO: Compare candidates using physically and project-relevant criteria.
% Do not select an approach solely on popularity or implementation ease.

\subsection{Assumptions and Applicability}
\label{sec:res-geo-spec-assumptions}

% TODO: State assumptions and the conditions under which they are valid.

\subsection{Limitations and Failure Conditions}
\label{sec:res-geo-spec-limitations}

% TODO: State where the evidence, model or selected approach ceases to be
% reliable or applicable.

\subsection{Project-Specific Conclusions}
\label{sec:res-geo-spec-conclusions}

The integrated RES-GEO output now supports the immediate project
scope: Tohoku/local comparison and barrier-class hydrodynamic
comparison. It does not authorise geometry generation, structural
damage simulation or ML surrogate training.

\subsection{Selected Approach and Rationale}
\label{sec:res-geo-spec-selected-approach}

The active approach is descriptor-first fixed-rigid barrier-class
comparison. Existing-defence reconstruction may be added where source
geometry is available. Unrestricted generative geometry, topology
optimisation, surrogate-led geometry search and deformable/damage
geometry are deferred until hydrodynamic extraction and validation
are stable.

\subsection{Unresolved Decisions}
\label{sec:res-geo-spec-unresolved-decisions}

\begin{itemize}
  \item TODO: citable Kamaishi bay-mouth breakwater geometry and 2011
    condition source.
  \item Final corridor-relative placement sections and geometry bounds.
  \item Geometry file format and HDF5/structured schema details.
  \item Roughness/opening/porous-closure activation policy.
\end{itemize}

\subsection{Requirements and Handoffs}
\label{sec:res-geo-spec-requirements}

Software may begin the geometry-manifest MVP in parallel with
Research: parse geometry case records, generate corridor-relative
barrier patches, preserve patch identifiers, run clipping/meshability
checks and export structured geometry metadata for replay-boundary
and local 3D scaffolds. This is a data-handling handoff only, not
solver implementation.

\subsection{Deliverable Completion Record}
\label{sec:res-geo-spec-completion-record}

% TODO: Record whether the Deliverable is:
%
% - Not started
% - In progress
% - Draft complete
% - Reviewed
% - Accepted
%
% Acceptance requires:
%
% - the research questions to be answered;
% - claims to be supported by appropriate evidence;
% - assumptions and limitations to be explicit;
% - the project-specific conclusion to be stated;
% - downstream requirements to be traceable;
% - unresolved decisions to be closed or formally deferred.


\clearpage
\printbibliography[
  heading=bibintoc,
  title={References}
]

\end{document}
